% Generated draft: requires mathematical and source review.
\documentclass[11pt]{article}
\usepackage[margin=25mm]{geometry}
\usepackage{fontspec}
\setmainfont{Latin Modern Roman}
\usepackage{amsmath,amssymb,mathtools,unicode-math}
\setmathfont{Latin Modern Math}
\usepackage{enumitem,needspace,xurl,hyperref}
\setlength{\parindent}{0pt}
\setlength{\parskip}{5pt}
\setlength{\emergencystretch}{3em}
\newcommand{\N}{\mathbb{N}}
\newcommand{\Z}{\mathbb{Z}}
\newcommand{\Q}{\mathbb{Q}}
\newcommand{\R}{\mathbb{R}}
\newcommand{\C}{\mathbb{C}}
\newcommand{\F}{\mathbb{F}}
\newcommand{\A}{\mathbb{A}}
\newcommand{\PP}{\mathbb{P}}
\newcommand{\E}{\mathbb{E}}
\newcommand{\eps}{\varepsilon}
\newcommand{\dd}{\,\mathrm{d}}
\begin{document}
\section*{UnsolvedMath ID: 20003162}
\textbf{Source problem code:} catalogue entry
\subsection{Definitions and mathematical statement}
A group action has cohomogeneity one when its principal orbits have codimension one. An Einstein metric satisfies $\operatorname{Ric}_g=\Lambda g$ for a constant $\Lambda$. The original problem asks to study existence of Einstein metrics on cohomogeneity-one manifolds; it does not specify manifold, action, completeness, Einstein sign or an extra sectional-curvature condition.

The aid treats the fixed standard action of $\operatorname{SO}(p+1)\times\operatorname{SO}(q+1)$ on $S^{p+q+1}$, $p,q\geq2$. On the regular part an invariant metric has the labeled two-warp form
\[
g=\dd t^2+a(t)^2g_{S^p}+b(t)^2g_{S^q},\qquad a,b>0,
\]
with unit round factor metrics. Its Einstein equations include
\[
a''=\frac{(p-1)(1-(a')^2)}a-q\frac{a'b'}b-\Lambda a,
\qquad
b''=\frac{(q-1)(1-(b')^2)}b-p\frac{a'b'}a-\Lambda b,
\]
and the radial equation $-pa''/a-qb''/b=\Lambda$.

A principal orbit is totally geodesic exactly when $a'=b'=0$ there, a stronger requirement than vanishing mean curvature. At such a regular time $t_0$, the smooth autonomous system and Cauchy uniqueness force
\[
a(t_0+s)=a(t_0-s),\qquad b(t_0+s)=b(t_0-s)
\]
throughout the corresponding regular interval. The aid's obstruction is that the standard sphere join diagram collapses the first labeled factor at one endpoint and the second at the other, incompatible with this identity-labeled reflection. Therefore no invariant Einstein metric for this fixed action has a totally geodesic principal orbit. This does not deny the existence of invariant Einstein metrics themselves, and even $p=q$ does not remove the labeled-Cauchy obstruction by exchanging factors.

For a different, same-collapse double diagram, reflection is compatible with the endpoint labels. For positive Einstein constant, scale to $\Lambda=p+q$. At a totally geodesic center, write $A=a(t_0)>0$, $B=b(t_0)>0$. The Hamiltonian constraint becomes
\[
\frac{p(p-1)}{A^2}+\frac{q(q-1)}{B^2}
=(p+q)(p+q-1).
\]
Thus the scale-normalized center data lie on one explicit one-parameter curve. Smooth closure still requires successful shooting to the designated singular orbit.

For a cap collapsing the $S^p$ factor, use inward distance $t$ from that cap. The smooth round-factor conditions are that $a$ extends oddly with $a'(0)=1$, while $b$ extends evenly with $b(0)=B_0>0$. In particular
\[
a=t+\alpha t^3+O(t^5),\qquad b=B_0+\beta t^2+O(t^4),
\]
with the Einstein Taylor diagnostics
\[
\beta=\frac{(q-1)/B_0-\Lambda B_0}{2(p+1)},\qquad
\alpha=-\frac{\Lambda+2q\beta/B_0}{6p}.
\]
The opposite cap uses reversed inward distance; the other collapse type interchanges labels. Parity and these initial coefficients are necessary diagnostics, not a claim that every point on the constraint curve reaches a smooth cap. This is an obstruction and a shooting reduction within a special action family, not a general Einstein-existence classification.
\subsection{Short English statement}
The supplied aid rules out a totally geodesic principal orbit for invariant Einstein metrics in the standard sphere join action. In a compatible same-collapse diagram it reduces positive-Einstein center data to one constraint curve, with cap parity and Taylor conditions still needed for shooting.
\subsection{Sources}
\begin{itemize}
\item Local UnsolvedMath registry, permanent ID $20003162$, source code \texttt{catalogue entry}; canonical \texttt{aim-topology-notes.json}, \texttt{notes[249]}.
\item AIM workshop \emph{Manifolds with nonnegative sectional curvature}, Problem 37: \url{https://aimath.org/WWN/nnsectcurvature/nnsectcurvature.pdf}. Original remarks and literature field are empty.
\item Research attempt 1: \texttt{partial\_result}, status at run \texttt{open}, current status \texttt{partially\_solved}, classification \texttt{PARTIAL-PROGRESS}, rigor \texttt{proved}, novelty confidence low. No authoritative complete literature reference is supplied for the proposed ODE obstruction and shooting reduction. The machine-generated aid is not peer reviewed and requires independent verification before citation.
\end{itemize}
\end{document}
